\question{10}{Considere un grafo representado por la siguiente matriz de adyacencia:

\begin{center}
\(
\begin{pmatrix}
0 & 1 & 0 & 1 & 0 & 0 \\
1 & 0 & 1 & 0 & 0 & 0 \\
0 & 1 & 0 & 0 & 1 & 0 \\
1 & 0 & 0 & 0 & 1 & 1 \\
0 & 0 & 1 & 1 & 0 & 1 \\
0 & 0 & 0 & 1 & 1 & 0
\end{pmatrix}
\)
\end{center}

La primera fila representa el vértice \(v_1\) y la última \(v_6\). Cada vez que visite los vecinos de un vértice, hágalo en orden ascendente de índice.
}

\begin{enumerate}
  \item (2 décimas) Dibuje el grafo que representa la matriz.
  \item (1 décimas) ¿El grafo es dirigido?
  \item (1 décimas) ¿El grafo es ponderado?
  \item (6 décimas) Ejecute BFS desde \(v_3\). Para cada paso (cada vez que desencole un vértice), muestre el contenido de la cola y el arreglo de predecesores en ese momento.
\end{enumerate}

\bigskip
